% Options for packages loaded elsewhere
\PassOptionsToPackage{unicode}{hyperref}
\PassOptionsToPackage{hyphens}{url}
\documentclass[
  11pt,
  a4paper,
]{article}
\usepackage{xcolor}
\usepackage[margin=1in]{geometry}
\usepackage{amsmath,amssymb}
\setcounter{secnumdepth}{-\maxdimen} % remove section numbering
\usepackage{iftex}
\ifPDFTeX
  \usepackage[T1]{fontenc}
  \usepackage[utf8]{inputenc}
  \usepackage{textcomp} % provide euro and other symbols
\else % if luatex or xetex
  \usepackage{unicode-math} % this also loads fontspec
  \defaultfontfeatures{Scale=MatchLowercase}
  \defaultfontfeatures[\rmfamily]{Ligatures=TeX,Scale=1}
\fi
\usepackage{lmodern}
\ifPDFTeX\else
  % xetex/luatex font selection
\fi
% Use upquote if available, for straight quotes in verbatim environments
\IfFileExists{upquote.sty}{\usepackage{upquote}}{}
\IfFileExists{microtype.sty}{% use microtype if available
  \usepackage[]{microtype}
  \UseMicrotypeSet[protrusion]{basicmath} % disable protrusion for tt fonts
}{}
\makeatletter
\@ifundefined{KOMAClassName}{% if non-KOMA class
  \IfFileExists{parskip.sty}{%
    \usepackage{parskip}
  }{% else
    \setlength{\parindent}{0pt}
    \setlength{\parskip}{6pt plus 2pt minus 1pt}}
}{% if KOMA class
  \KOMAoptions{parskip=half}}
\makeatother
\usepackage{longtable,booktabs,array}
\newcounter{none} % for unnumbered tables
\usepackage{calc} % for calculating minipage widths
% Correct order of tables after \paragraph or \subparagraph
\usepackage{etoolbox}
\makeatletter
\patchcmd\longtable{\par}{\if@noskipsec\mbox{}\fi\par}{}{}
\makeatother
% Allow footnotes in longtable head/foot
\IfFileExists{footnotehyper.sty}{\usepackage{footnotehyper}}{\usepackage{footnote}}
\makesavenoteenv{longtable}
\setlength{\emergencystretch}{3em} % prevent overfull lines
\providecommand{\tightlist}{%
  \setlength{\itemsep}{0pt}\setlength{\parskip}{0pt}}
\usepackage{bookmark}
\IfFileExists{xurl.sty}{\usepackage{xurl}}{} % add URL line breaks if available
\urlstyle{same}
\hypersetup{
  hidelinks,
  pdfcreator={LaTeX via pandoc}}

\author{}
\date{}

\begin{document}

{
\setcounter{tocdepth}{3}
\tableofcontents
}
\section{Controfase as Phase-Inversion Operator on the Reactive Circuit:
An Ordinative Reading of Standing Waves in a Spinning
Fluid}\label{controfase-as-phase-inversion-operator-on-the-reactive-circuit-an-ordinative-reading-of-standing-waves-in-a-spinning-fluid}

\subsection{A working paper on Singh, Rønning, Liu, Angheluta, Concha,
and Bandi (Communications Physics, 2026), with formalization of the
phenomenological signature of Controfase and consequences for the
ordinative sciences research
programme}\label{a-working-paper-on-singh-ruxf8nning-liu-angheluta-concha-and-bandi-communications-physics-2026-with-formalization-of-the-phenomenological-signature-of-controfase-and-consequences-for-the-ordinative-sciences-research-programme}

\textbf{Author}: Fabio Ghioni \textbf{ORCID}:
\href{https://orcid.org/0009-0009-0415-9434}{0009-0009-0415-9434}
\textbf{Affiliation}: Ordinative Sciences Foundation ETS ---
president@ordinativescience.foundation \textbf{Framework versions}:
TE\_CORE v5.1 / OST v2.1 / OST Extension (Teleodynamics \& Causal
Inversion) v1.1 \textbf{Document status}: Working paper. Constitutes (i)
a structural reading of a peer-reviewed experimental result within the
Ordinative Sciences framework, (ii) formalisation of the
Phenomenological Signature of Controfase as cross-domain diagnostic
instrument, (iii) foundational material for the volume on Controfase,
the Reaction-Diffusion Civilizational Dynamics paper v1.3, the
Ordinative Category Theory programme, and adjacent work in the
Ordinative Sciences corpus. \textbf{Version}: 2.1 \textbf{Date}: April
23, 2026 \textbf{Metadata revision}: October 2026 --- institutional
metadata updated to the Ordinative Sciences Foundation ETS; text
unchanged \textbf{License}: Creative Commons Attribution 4.0
International (CC BY 4.0) \textbf{DOI}:
\href{https://doi.org/10.5281/zenodo.23439941}{10.5281/zenodo.23439941}

\begin{center}\rule{0.5\linewidth}{0.5pt}\end{center}

\subsection{Abstract}\label{abstract}

A recent experimental and theoretical paper by Singh, Rønning, Liu,
Angheluta, Concha, and Bandi (\emph{Communications Physics} 9:123, 2026)
reports that shallow-water standing waves scattered by a single
irrotational vortex generate quantized, system-spanning nodal lines
rotating counter to the vortex circulation. The result generalizes a
classical Aharonov--Bohm hydrodynamic analogue due to Berry et
al.~(1980): where the original Berry configuration produces wavefront
dislocations localized near the vortex core, the standing-wave
configuration produces global topological structures whose number is
fixed by a single dimensionless parameter derived from a continuous
source. The authors document an unexpected fact: their prior theoretical
intuition predicted either cancellation or superposition; neither
occurred. The first appearance of the nodal lines was attributed to
experimental artefact and only accepted as a genuine phenomenon after
independent reproduction in numerical simulations.

This paper offers an ordinative reading of the Singh et al.~result. The
reading identifies the experimental configuration as an instance of
Controfase --- defined within the Technology of Expressions framework
(Ghioni, TE\_CORE §7) as a universal ordinative operator that introduces
a phase-shift in the automatic stimulus--response sequence of a reactive
system, interrupts inertia, and reopens the coherence field. Four
propositions are developed: (i) the Berry-to-Singh transition is a
phase-inversion of the reactive circuit between attractor (vortex) and
expressive terminal (wave field), not a reactive compensation; (ii) the
energy emitted by the attractor is redirected to the relational field
\(R\) and used to invert the passive--active polarity of the terminals,
with no energetic cost to the recipient system; (iii) the researchers'
own reported sequence of reactions --- radical disconfirmation of prior
intuition, initial dismissal of the observation as artefact, abandonment
of the working program to pursue the phenomenon --- constitutes a
phenomenological signature of Controfase that is operationally
diagnostic and can be formalized as a three-marker criterion; (iv) the
hydrodynamic case therefore furnishes both a structural exemplar of the
Controfase operator in a classical physical system and an external,
peer-reviewed, reproducible empirical anchor for the ordinative claim
that Controfase operates across the full hierarchy of configured
systems, from the quantum and classical-physical through the biological
to the cognitive and social.

The paper develops the ontological framework required to make these
propositions intelligible to readers approaching the ordinative sciences
for the first time, formalizes the three-marker phenomenological
signature as a diagnostic criterion applicable across domains, and
records consequences for the ordinative sciences research programme.
Confidence grades (\(S_0\)--\(S_3\)) are maintained throughout. An
appendix provides a glossary of ordinative constructs.

\begin{center}\rule{0.5\linewidth}{0.5pt}\end{center}

\subsection{1. Introduction}\label{introduction}

\subsubsection{1.1 The experimental
result}\label{the-experimental-result}

The Aharonov--Bohm (AB) effect, predicted in 1959, showed that quantum
observables depend on electromagnetic potentials and not only on local
fields: a charged particle acquires a phase shift on encircling a region
of confined magnetic flux even when the particle itself traverses no
region of nonzero field (Aharonov \& Bohm, 1959). In 1980 Berry,
Chambers, Large, Upstill, and Walmsley demonstrated that AB has a
precise classical analogue in shallow-water surface waves scattered by
an irrotational draining vortex, with the vector potential replaced by
the azimuthal velocity field of the vortex, the magnetic flux replaced
by the circulation \(\Gamma\), and the scattering parameter given by

\[ \alpha_q = \frac{q\Phi_B}{h} \;\longleftrightarrow\; \alpha_f = -\frac{\Gamma}{\lambda v_g} \tag{1.1} \]

In the Berry configuration a single travelling wave develops wavefront
dislocations --- localized topological defects pinned near the vortex
core, whose number is set by \(\alpha\).

Singh, Rønning, Liu, Angheluta, Concha, and Bandi (2026) modify exactly
one structural parameter of the Berry setup. Instead of a single
travelling wave, they generate a standing wave by superposing two
counter-propagating travelling waves of equal wavelength, with the
vortex at the centre. The scattering problem acquires double angular
periodicity, expressed mathematically by the substitution \(n \to 2n\)
in the Bessel-function summation:

\[ \psi = 2\sum_{n} (-i)^{l_q} J_{l_q}(kr) \, e^{2 n i \varphi}, \qquad l_q = |2n - \alpha| \tag{1.2} \]

The empirical and theoretical consequence is qualitatively new. Instead
of localized wavefront dislocations near the vortex core, the system
develops a discrete set of \emph{nodal lines} --- system-spanning lines
of zero wave amplitude --- that radiate outward, rotate with angular
velocity \(\omega = \nu / \alpha\), and rotate in the direction
\emph{opposite} to the vortex. The number \(N\) of nodal lines is
bounded by

\[ |\alpha| - 1 < N < |\alpha| + 1 \tag{1.3} \]

Integer \(\alpha\) yields exactly \(|\alpha|\) lines; non-integer
\(\alpha\) yields oscillation between the two integers bracketing
\(|\alpha|\). The quantization is \emph{topological} and emerges from
the geometry of the configuration, not from any quantization of the
source: \(\Gamma\) is continuous and freely tunable. The authors
validate the prediction with high-speed caustic imaging, particle
imaging velocimetry, and analytical asymptotics, with agreement between
experiment and theory reported as ``no adjustment parameters.'' The
paper is open access, peer-reviewed, and the analysis code is archived
(Zenodo DOI 10.5281/zenodo.18652566).

\subsubsection{1.2 The unexpected character of the
result}\label{the-unexpected-character-of-the-result}

The Singh et al.~paper contains four reported facts that, in the present
reading, are not incidental. They are the phenomenological signature the
experiment carries, and they are the starting point of the ordinative
analysis that follows.

First, Rønning (co-first author) reports that the team expected either
cancellation or superposition of the dislocation patterns, and states
that ``our intuition was completely wrong.'' The language is not
neutral: ``completely wrong'' registers a radical disconfirmation, not
an incremental refinement. The theoretical framework the researchers
brought to the experiment did not contain the observed response as a
possibility.

Second, Singh (co-first author) reports that when the nodal lines first
appeared, the team interpreted them as experimental artefact. The
phenomenon was not accepted as genuine on first observation; it was
rejected as spurious. Only when independent numerical simulations
produced the same structures was the observation promoted from artefact
to fact.

Third, Singh further reports: ``we dropped everything and quickly worked
out the mathematics underlying how they arise.'' The phenomenon
reorganized the research programme; the analytical framework was
constructed \emph{after} the observation, to chase what the experiment
had already shown.

Fourth, Bandi (senior author), commenting on the planned extension to
vortex lattices emulating superconductor vortex matter, says: ``we don't
yet know what we'll see --- and that's exactly what makes it worth
doing.'' This is not a formulaic disclaimer; it is an explicit
acknowledgement that predictive theory does not anticipate the system's
response in the next configuration.

These four reports together describe a specific phenomenological
pattern. A reactive epistemic system (the researchers' prior theoretical
framework, functioning as the ``attractor'' of prediction) emits an
expectation (cancellation or superposition). The system being probed
(the fluid configuration) produces a response that does not complete the
predicted reactive circuit. The researchers, whose attention is
structured by the predictive framework, first interpret the response as
noise (artefact); on verification they abandon the framework (dropped
everything) and reorganize around the phenomenon. The character of the
encounter is surprise in a structural, not rhetorical, sense: the
system's response was not contained in the predictive framework that was
probing it.

\subsubsection{1.3 The ordinative reading proposed in this
paper}\label{the-ordinative-reading-proposed-in-this-paper}

The thesis developed in what follows is that the Singh et
al.~configuration realizes the structural pattern that the Technology of
Expressions (TE) framework has formalized under the name
\textbf{Controfase}. The Italian term, introduced by the first author of
this paper in earlier TE literature and retained without translation for
citation continuity, denotes a universal ordinative operator that
introduces a phase-shift in the automatic stimulus--response sequence of
a reactive system, interrupting the system's inertia and reopening the
coherence field. The term has been developed primarily in contexts of
cognitive and social dynamics, where its operative features are well
documented; this paper contributes the structural reading of a
controlled physical system in which the same operator is argued to be at
work.

Four specific propositions are advanced:

\textbf{Proposition 1} (Structural). The transition from the Berry
single-wave configuration to the Singh standing-wave configuration is a
phase-inversion of the reactive circuit between the vortex (attractor)
and the wave field (expressive terminal), not a reactive compensation.
The nodal lines are not a pattern of resistance or cancellation; they
are the geometry produced when the reactive energy emitted by the vortex
is intercepted at the circuit-level, inverted in phase, and redirected
through the relational structure of the configured system.

\textbf{Proposition 2} (Energetic). The energy that would, in a purely
reactive system, produce action--reaction scattering near the vortex
core is redirected to the relational field \(R\) and used to invert the
passive--active polarity of the terminals (vortex ↔ field). The vortex
finances --- energetically --- the configuration in which its own
organising influence produces a global structure rotating against it. No
energy cost is borne by the wave field as a separate reservoir: the
field's coherent response is constructed with the attractor's own
energy, phase-inverted.

\textbf{Proposition 3} (Phenomenological). The researchers' reported
sequence of reactions constitutes an operationally diagnostic signature
of Controfase. Three markers are identifiable: (a) radical
disconfirmation of a prior intuitive or theoretical framework, (b)
initial interpretation of the anomalous observation as artefact, (c)
abandonment of the active research programme to reorganize around the
phenomenon. This three-marker signature is generalizable and
formalizable as a diagnostic criterion for the presence of Controfase in
other domains.

\textbf{Proposition 4} (Ontological). The hydrodynamic case instantiates
Controfase in a classical physical system without cognitive or social
mediation. This establishes that Controfase is a universal operator
acting across scales --- from quantum and classical-physical
configurations through biological and cognitive systems to social and
cultural ones. Intelligence in the ordinative sense --- structural
adaptivity toward dynamic equilibrium, distinct from cognitive agency
--- is a graded property of configured matter, present throughout the
hierarchy of configured systems, and fully expressed in systems with
sufficient structural coherence to host the Controfase operator stably.

The paper develops these four propositions, establishes their confidence
grades individually, formalizes the phenomenological signature as a
diagnostic criterion, and records consequences for the active ordinative
sciences research programme. The final section addresses the analytical
perimeter of the work.

\begin{center}\rule{0.5\linewidth}{0.5pt}\end{center}

\subsection{2. The experimental result, restated
structurally}\label{the-experimental-result-restated-structurally}

To make the ordinative reading tractable, the Singh et al.~experiment is
here restated in a form that separates physical content from
configuration. This reformulation does not alter the paper's physics; it
isolates the structural features on which the ordinative analysis is
carried out.

\textbf{Constituents}. Three classes of element are co-present in the
experimental domain:

\begin{itemize}
\tightlist
\item
  A \textbf{stationary irrotational vortex} at the centre of the
  apparatus, maintained by a constant drain rate. Far-field velocity
  field \(\vec U(\vec r) = (\Gamma / 2\pi r)\hat\varphi\), with
  circulation \(\Gamma\) continuously tunable by the experimenter.
\item
  A \textbf{surface wave field} generated by two pairs of acoustic
  transducers vibrating in phase on opposite sides of the vortex,
  producing two counter-propagating travelling waves of common
  wavelength \(\lambda\). Their superposition constitutes a standing
  wave with the vortex at its centre.
\item
  The \textbf{fluid medium} of average depth \(H\), satisfying the
  shallow-water equation
  \[ (\partial_t + \vec U\cdot\nabla)^2 \eta - c^2 \nabla^2\eta = 0, \qquad c = \sqrt{gH}. \tag{2.1} \]
\end{itemize}

\textbf{Configuration regime}. The experiment operates in the
shallow-water short-wavelength weak-flow regime: \(hk \ll 1\),
\(U \ll c\), \(\beta = Rk \gg 1\). In this regime the surface
deformation \(\eta\) maps onto a Schrödinger equation for a charged
particle in the azimuthal gauge of the vortex flow, and the scattering
problem admits the analytical treatment given by Equations (1.2) and
(1.3).

\textbf{Observed response}. The standing-wave-plus-vortex system
develops a discrete set of nodal lines whose properties are (i) globally
system-spanning, subject to dissipative attenuation in the far field,
(ii) integer-quantized in number, (iii) counter-rotating with respect to
the vortex at angular velocity \(\omega = \nu / \alpha\), (iv)
phase-opposed on either side, with a \(\pi\) phase shift of \(\eta\)
across each line, and (v) independent of the near-core micro-structure
of the vortex, depending only on the far-field background flow.

\textbf{Structural contrast with the Berry configuration}. The critical
observation for the ordinative reading is that the transition from the
Berry configuration (\(S_B\)) to the Singh configuration (\(S_S\))
modifies exactly one structural feature: the relational structure \(R\)
that specifies how the wave field is incident on the vortex. In \(R_B\)
the incidence is single-direction; in \(R_S\) it is symmetric
double-direction. The constituents \(\Sigma_B\) and \(\Sigma_S\) differ
by the addition of one travelling wave; the governing equations are
unchanged; the medium is unchanged; the vortex is unchanged. And yet the
response \(\Phi\) transitions from localized wavefront dislocations to
globally structured counter-rotating nodal lines. This is a qualitative
change in response produced by a structurally minimal change in
configuration.

The authors observe this contrast and treat it, correctly within their
framework, as an instance of the sensitivity of topological scattering
solutions to boundary and symmetry conditions. The ordinative reading
does not displace this observation; it asks what \emph{kind} of
structural pattern is being instantiated, across its appearance in
different substrates, when a minimal change in configuration produces a
qualitative change in response.

\begin{center}\rule{0.5\linewidth}{0.5pt}\end{center}

\subsection{3. Ontological framework: structural intelligence and the
Controfase
operator}\label{ontological-framework-structural-intelligence-and-the-controfase-operator}

This section establishes the ontological framework within which the
readings of §§ 4--8 operate. It is directed at readers approaching the
ordinative sciences for the first time. The constructs developed here
are foundational; the analysis that follows assumes them. Readers may
test the constructs against their own categories and against the
empirical evidence offered in §§ 4--8 --- which is the appropriate way
to engage a foundational proposal --- but the argument of this paper
does not retreat from these constructs to softer alternatives.

\subsubsection{3.1 Reactive systems and the third
law}\label{reactive-systems-and-the-third-law}

A \emph{reactive system}, in the usage developed here, is one whose
response to an action is determined by the structure of the action
impinging on it, independently of any internal organization of the
system toward dynamic equilibrium. Newton's third law --- every action
produces an equal and opposite reaction --- is the formal statement of
the response structure of an ideal reactive system. The incident force
encounters a structure whose character is to resist; the resistance
generates the return force; the event-interaction concludes with an
energetic exchange that is symmetric in magnitude, localized in space
and time, and predictable from the properties of the incident action and
the resistive structure.

The ideal reactive system is an \emph{idealization}. Real systems
exhibit reactive behaviour in varying degrees; very few systems are
purely reactive in the strict Newtonian sense. A rigid steel wall
approximates the ideal in its response to a low-energy mechanical impact
on timescales short compared to the wall's own internal dynamics. A gas
at equilibrium is less purely reactive than the wall, because its
molecular structure allows redistribution of the incident energy across
internal degrees of freedom. A biological cell, under perturbation, is
less purely reactive still, because its regulatory networks permit
adaptive response on timescales comparable to the perturbation. A human
being is the least purely reactive: the same stimulus can produce
radically different responses depending on internal configuration,
attentional state, identity structure, and context.

The gradient from the steel wall to the human being is not a gradient in
\emph{complexity} in the conventional sense, though complexity varies
along it. It is a gradient in \textbf{structural intelligence} as
understood in the TE framework. Structural intelligence, in this usage,
is the capacity of a configured system to respond to a stimulus in a way
that is oriented toward the system's dynamic equilibrium (its own, or
that of the larger system of which it is part) rather than toward the
completion of the Newtonian reactive circuit. Intelligence in this sense
is not cognitive agency; it is a structural property of configured
systems that admits degrees.

The commitment is that structural intelligence is present, in attenuated
form, throughout configured matter. Quantum systems exhibit it in the
contextuality of measurement outcomes, which is not reducible to
reactive dynamics (Bell, 1964; Kochen \& Specker, 1967). Classical
systems exhibit it in the adaptive character of dissipative structures
(Prigogine \& Stengers, 1984) and in the non-additive responses of
configured fluids, solids, and fields. Living systems exhibit it in
regulatory and homeostatic dynamics. Cognitive systems exhibit it in
attention, choice, and reflective reorganization. The TE framework
treats these as manifestations along a single structural axis, varying
in depth and stability of expression, not as categorically distinct
phenomena.

This is a foundational position of the TE framework, entailed by Axiom 0
(universality under isomorphism and synesthesia; TE\_CORE §2.0) and
Axiom 14 (the Ordinative Set as relational field; TE\_CORE §2.5). The
present paper articulates the position explicitly because the analysis
of §§ 4--8 is conducted within it.

\subsubsection{3.2 The Controfase
operator}\label{the-controfase-operator}

Within a reactive system, the completion of the stimulus--response
circuit is automatic. A force incident on a rigid wall produces a
reaction force; an insult directed at a person in reactive configuration
produces anger, withdrawal, or counter-attack; a legal threat directed
at a person in reactive configuration produces anxiety, defensive
preparation, or collapse. The reactive response is predictable from the
stimulus and the structural configuration of the respondent.

The \textbf{Controfase operator} (TE\_CORE §7) is the universal
ordinative operator that, when active in a sufficiently structured
system, intercepts the stimulus--response circuit \emph{before its
automatic completion}, phase-inverts the reactive impulse, and redirects
the energy of the incident stimulus to invert the passive--active
polarity between stimulus-source and respondent. The stimulus does not
meet the expected reactive structure; the energy of the stimulus is not
returned to the source as resistance or scattered in dissipation; it is
used to reconfigure the relational field such that the anticipated
event-interaction does not occur.

The canonical operational example in cognitive-social contexts is the
response to an insult that takes the form of full agreement with the
insult's content, stated without irony, without strategic calculation,
and without internal reactive discomfort: \emph{``You are right, I am an
idiot who will never amount to anything.''} The insult was emitted with
the intention of producing a defensive reactive response. The response,
made on the same vector as the insult, deprives the insult of its
expected resistive structure. The insulter finds no resistance to push
against; the energy expended in emitting the insult is not returned as
resistance; no event-interaction of the expected form takes place. The
insulter is left with their own emitted energy, now unattached to a
reactive target, and frequently reports subjective disorientation.

This response is not tactical. A tactical response involves strategic
reading of the insulter's intent and deliberate choice to adopt a
non-reactive posture as a manipulative move. Tactical responses break
down under sustained pressure because the underlying reactive circuit
remains active and the suppression is performed against it. The
Controfase response is structural: the person in Controfase
configuration has an internal organization in which the insult's content
does not find an automatic reactive circuit to activate. The response is
produced by the same internal configuration that made the automatic
reactive response unavailable; it is not a performance overlaid on a
suppressed reaction. The distinction is decisive: the structural
Controfase is stable across contexts, the tactical performance is not.

A second canonical example is the response to a formal legal threat (a
summons, a denunciation, a suit). The reactive response is anxiety,
activation of defensive systems, cognitive preoccupation. The Controfase
response associates the event with satisfaction, fulfilment, or a sense
of completion --- \emph{without} cognitive reframing or strategic
calculation --- because the internal configuration does not organize the
event as a threat. The legal action, regardless of its formal character,
does not meet the reactive target the issuer's framework presupposed.

\subsubsection{3.3 The energetic
accounting}\label{the-energetic-accounting}

The Controfase operator does not violate conservation of energy. Energy
emitted by the stimulus-source is emitted and must be accounted for. The
question is where it goes.

In the reactive case the accounting is closed at the event-interaction:
incident kinetic energy becomes rebound kinetic energy plus deformation,
heat, or sound in the interacting bodies. The total is conserved; the
trajectories of the interacting bodies are altered by the exchange.

In the Controfase case, the accounting is more subtle. The energy
emitted by the stimulus-source is not received by the respondent as
scattering energy; it is received by the \emph{relational field} \(R\)
between source and respondent. \(R\) is reconfigured. The
reconfiguration of \(R\) is what produces the response --- the response
that inverts the passive--active polarity. The energy spent by the
stimulus-source is therefore used to construct the configuration in
which the expected reactive event does not take place and the source's
position in the polarity has been inverted.

The respondent bears no energetic cost in the strict sense: the
respondent does not expend energy proportional to the stimulus in order
to produce the response. The response is constructed with the stimulus's
own energy, phase-inverted. The respondent's cost, in Controfase, is the
\emph{structural cost} of maintaining the internal configuration that
makes the Controfase response available --- a cost that is independent
of the particular stimulus, paid before the stimulus arrives, and
amortized across all stimuli that encounter the configured respondent.
This structural cost is the prior stabilization of the respondent's
coherent identity, in the TE usage (TE\_CORE §§ 4, 7).

The stimulus-source, by contrast, has spent energy to emit and has
received nothing in return that it can use to complete its intended
action. In the paradigmatic examples, this manifests as the
residual-energy phenomena documented in TE\_MODULE\_PPRO: the insulter
remains agitated after the interaction; the issuer of the legal action
does not obtain the expected destabilization; the attacker's own energy
remains unresolved in its internal system. The Controfase does not
reflect the energy back at the source (as a mirror response would); the
Controfase redirects the energy to reconfigure \(R\), leaving the
source's own internal imbalance intact.

\subsubsection{3.4 Structural intelligence as a
spectrum}\label{structural-intelligence-as-a-spectrum}

This framing has an immediate consequence. Controfase is a universal
operator available to all sufficiently structured systems; structural
intelligence is graded across configured matter. Controfase is therefore
observable in systems across the full spectrum --- physical, chemical,
biological, informational, cognitive, social --- in forms appropriate to
each substrate's structural complexity.

The articulation of Controfase in TE literature has so far emphasized
the cognitive-social domain, where its operational features are most
salient and where its pathological absence (the entropic engram;
TE\_CORE §§ 4, 5.4) is most destructive to the systems that exhibit it.
Frontier applications of the operator to physical phenomena have been
developed in the framework's working notes --- gravity reconstrued as
permanent Controfase against entropic dispersion; magnetism as polarized
Controfase; temporal dilation as a function of a system's Controfase
frequency; non-Newtonian propulsion as Controfase on the spatial
relational field --- and the present paper anchors a specific class of
hydrodynamic Controfase to a controlled, peer-reviewed, reproducible
experimental result.

The following four sections develop the analysis.

\begin{center}\rule{0.5\linewidth}{0.5pt}\end{center}

\subsection{4. Reading I: Controfase as phase-inversion operator in the
Singh--Bandi
configuration}\label{reading-i-controfase-as-phase-inversion-operator-in-the-singhbandi-configuration}

\subsubsection{4.1 The reactive circuit of the Berry
configuration}\label{the-reactive-circuit-of-the-berry-configuration}

In the Berry (1980) configuration, a single travelling wave is incident
on the vortex. The circuit structure is the following. The wave carries
momentum in the direction of propagation. The vortex, through its
azimuthal velocity field, produces a phase shift on the portions of the
wavefront passing on either side of its core. Where the phase shifts do
not match, wavefront dislocations form: topological defects in the
wavefront at which the phase is undefined. The dislocations are
\emph{localized} in a neighbourhood of the vortex core. The wave,
downstream, recovers its structure with a phase offset.

This is the structure of a reactive circuit in the sense developed in
§3. The wave acts on the vortex (the vortex scatters the wave); the
vortex acts on the wave (the wave acquires dislocations); the
event-interaction is localized, the energetic accounting is closed in a
neighbourhood of the vortex, and the downstream behaviour is determined
by the incoming behaviour modulo a calculable phase offset. The analogy
with Newtonian mechanics is structural, not merely rhetorical: the
vortex plays the role of the obstacle, the wave the role of the
impinging force, and the dislocations the role of the deformation
produced by the interaction. The reactive circuit closes.

\subsubsection{4.2 The interception and inversion in the Singh
configuration}\label{the-interception-and-inversion-in-the-singh-configuration}

The Singh standing-wave configuration modifies the circuit at the
structural level. Two waves, counter-propagating, are incident on the
vortex from opposite sides. The naïve expectation --- reported by the
authors --- is that the two dislocation patterns superpose, or that they
cancel by symmetry. Neither occurs. Instead the field adopts a global
configuration in which the vortex's organising influence is expressed
not as localized dislocations but as \emph{system-spanning nodal lines}
rotating against the vortex.

The ordinative reading of this transition is that the standing-wave
configuration intercepts the reactive circuit of the single-wave
configuration before its automatic completion and phase-inverts the
relational structure. In the single-wave case, the vortex and the wave
field stand in a source-receiver polarity: the vortex is active (the
source of the azimuthal flow), the wave is passive (the recipient of the
dislocation). In the standing-wave case, this polarity is inverted. The
wave field, by adopting its global nodal-line geometry, becomes the
active structural element expressing the vortex's influence: the vortex
continues to rotate, but the \emph{expression} of its rotation is now
the counter-rotating structure of the field, not the localized defects
of its direct scattering. The vortex is, structurally, reconfigured as
the organizing element whose organising action is carried by the field's
response, not as the scatterer of a wave's dislocation pattern.

This is the defining structure of Controfase. The vortex's reactive
emission --- the azimuthal velocity field that, in the Berry
configuration, produces a local scattering event --- is intercepted by
the double-periodicity of the standing wave before it can complete its
reactive event. The energy that would have gone into local scattering is
redirected through \(R\) (the relational structure of the standing-wave
configuration) and emerges as the global counter-rotating structure of
the nodal lines.

The following features of the observed response are the structural
signatures that confirm the reading.

\textbf{Signature A: Counter-rotation}. The nodal lines rotate with
angular velocity \(\omega = \nu/\alpha\), where the sign of \(\alpha\)
is opposite to the sign of \(\Gamma\). This is not a coincidental
counter-rotation; it is a counter-rotation whose rate is fixed by the
vortex's own circulation parameter with sign reversed. The field's
response \emph{uses} the vortex's parameters and \emph{inverts} them.
This is phase-inversion in the strict sense: same structural parameters,
inverted sign.

\textbf{Signature B: Quantization without source quantization}. The
number of nodal lines is integer-valued, with
\(|\alpha| - 1 < N < |\alpha| + 1\). The integer count is not inherited
from any quantized source --- \(\Gamma\) is continuous --- but is
generated by the topological constraint of the standing-wave
configuration acting on the continuous parameter. In ordinative terms,
this is the signature of \emph{structural quantization}: the relational
field \(R\) admits only discrete realizations that maintain global
coherence, even though the driving parameter is continuous. This is the
operative form, in classical fluid dynamics, of the integer-valued
structures the TE framework attributes to the action of Controfase on a
reactive field (TE\_CORE §7.3, formalization of the Controfase operator
\(C\) that decouples the automatic closure of the transition
\(s_{t+1} = f(s_t)\)).

\textbf{Signature C: Phase opposition across the nodal line}. The
surface deformation \(\eta\) acquires a \(\pi\) phase shift across each
nodal line, so that adjacent angular sectors of the field oscillate in
anti-phase. The field is not in a null state (it is not absent); it is
in a structured state whose adjacent zones are \emph{in counter-phase}.
This is the local signature of the global phase-inversion: across the
nodal line, the field instantiates the phase-displacement that defines
the Controfase operator at the global level.

\textbf{Signature D: Independence of attractor micro-structure}. Singh
et al.~note explicitly that the nodal line structure depends only on the
far-field circulation \(\Gamma\) and is ``robust to near-core details.''
This is the signature of an operator acting at the level of \(R\)
(relational structure) rather than at the level of \(\Sigma\)
(constituents): the specific internal dynamics of the vortex do not
enter the response. What matters is only the organising influence the
vortex exerts on the field as a whole. This is the expected behaviour of
Controfase, which operates on the reactive \emph{circuit} (the
stimulus--response relational structure) rather than on the internal
details of the stimulus-source.

\subsubsection{4.3 The energetic
accounting}\label{the-energetic-accounting-1}

The energy accounting for the Singh configuration confirms the reading.
The vortex is maintained by a constant drain rate; its energy input is
measurable and continuous. In the Berry single-wave configuration this
energy is partly returned to the wave (as the phase offset imprinted on
the downstream field) and partly dissipated viscously in the
neighbourhood of the core, where wavefront dislocations are established.
The reactive circuit closes locally.

In the Singh standing-wave configuration the same vortex, with the same
\(\Gamma\), sustains a qualitatively different field response. The
energy required to maintain the counter-rotating nodal-line structure
across the experimental domain is not paid by a separate energy source;
no additional actuator drives the nodal lines. The energy that sustains
the global structure is the same energy that the vortex continuously
provides through its azimuthal flow. The vortex finances, literally, the
configuration that expresses the vortex's organising influence as a
counter-rotation.

This is the structural feature that makes the reading precise. The wave
field does not expend energy of its own in resisting the vortex; the
wave field does not scatter the vortex's influence back as reactive
rebound; the wave field \emph{uses} the vortex's energy to construct the
geometry in which the expected reactive event does not take place. The
nodal lines rotate against the vortex using the vortex's own rotation as
their source, phase-inverted. No other energy flow is required; the
dissipative attenuation observed in the far field is secondary to the
primary energetic structure and represents the boundary of the
configuration's coherence in the presence of viscous losses.

\subsubsection{4.4 Confidence grade}\label{confidence-grade}

The identifications of §§ 4.1--4.3 are structural mappings of features
explicit in the Singh et al.~paper onto the formal definition of the
Controfase operator given in TE\_CORE §7. The features cited ---
counter-rotation with inverted sign of the attractor parameter, integer
quantization of the field response, \(\pi\) phase opposition across the
nodal locus, independence from attractor micro-structure, energetic
financing of the response by the attractor --- are all established in
the peer-reviewed paper at \(S_0\). The mapping of these features onto
the Controfase formalism is a direct correspondence. The reading of §§
4.1--4.3 is therefore rated \(S_1\).

The identification of the Controfase operator at work in the
hydrodynamic configuration with the operator at work in cognitive-social
configurations carries the universality proposition developed in §3 and
is rated \(S_2\) pending the further substrate-by-substrate
documentation that the ordinative research programme is producing.

\begin{center}\rule{0.5\linewidth}{0.5pt}\end{center}

\subsection{5. Reading II: the phenomenological signature of
Controfase}\label{reading-ii-the-phenomenological-signature-of-controfase}

\subsubsection{5.1 The observer's encounter with
Controfase}\label{the-observers-encounter-with-controfase}

The Singh et al.~paper documents not only the physical result but also,
in its published language, the sequence of reactions of the research
team to the result. Section 1.2 above recorded these reactions. The
ordinative reading proposed here is that the researchers' reactions
constitute a \emph{second} Controfase event, running in parallel with
the physical one, and exhibiting the same structural pattern at the
epistemic level. The formalization of this observation yields a
three-marker phenomenological signature that can be used as a diagnostic
criterion in other domains.

The predictive framework the researchers brought to the experiment is,
in ordinative terms, a reactive epistemic configuration. The framework
anticipated a specific class of responses: cancellation or superposition
of dislocation patterns, both of which would have closed the reactive
circuit between prediction and observation in the standard way. The
actual response --- nodal lines rotating against the vortex,
system-spanning, quantized --- did not close the reactive circuit. The
predictive framework therefore encountered the same structural pattern
that the single-wave configuration encountered when extended to standing
waves: a phenomenon that did not complete the expected reactive event.

The researchers' sequence of reported responses is the signature of this
epistemic Controfase encounter. Three markers appear with operational
clarity.

\subsubsection{5.2 Marker 1: Radical disconfirmation of prior
intuition}\label{marker-1-radical-disconfirmation-of-prior-intuition}

Rønning's report --- ``our intuition was completely wrong'' ---
registers a disconfirmation that is qualitatively different from
ordinary empirical updating. Ordinary empirical updating refines a
framework by adjusting its parameters in the direction of the data; the
framework survives, augmented or corrected. ``Completely wrong'' signals
that the framework did not contain the observed response as a
possibility and therefore cannot be adjusted to accommodate it without
restructuring. The phenomenon exceeds the space of possibilities the
framework is capable of generating.

This is the signature, at the epistemic level, of a stimulus that does
not find its expected reactive circuit. The framework --- the
``insulter'' in the structural analogy --- emitted an expectation and
the system did not return the anticipated reactive response. The
researchers recognize this as a radical, not incremental, failure of the
predictive structure.

In the cognitive-social domain, the parallel marker is the phenomenon
described in TE\_MODULE\_PPRO §9: a control attempt (manipulation,
coercion, social pressure) encounters a respondent whose internal
configuration does not actualize the reactive circuit the attempt
presupposed. The control attempt does not return to the source as the
expected dependency, acquiescence, or escalation. The structural
description is the same: radical mismatch between the expectation
carried by the stimulus and the response structurally available in the
respondent.

\subsubsection{5.3 Marker 2: Initial interpretation of the phenomenon as
artefact}\label{marker-2-initial-interpretation-of-the-phenomenon-as-artefact}

Singh's report --- the first observation of the nodal lines was
interpreted as experimental artefact --- is the characteristic
\emph{second response} of a reactive epistemic system to a Controfase
encounter. The mismatch between expectation and observation is first
interpreted as a defect in the observational apparatus rather than as a
feature of the observed system. The attribution of the anomaly to
artefact is the epistemic equivalent of a reactive system's attempt to
reject an out-of-circuit stimulus as noise: ``this should not have
happened; therefore it did not really happen; therefore the appearance
of its happening must be spurious.''

The attribution is not, in itself, a failure of scientific rigour; it is
the expected provisional response of an epistemic system to an
observation that does not fit its predictive framework. The rigour lies
in what is done next. The Singh et al.~team tested the observation
against independent numerical simulations; the simulations reproduced
the structures; the attribution to artefact was withdrawn; the
observation was promoted from noise to phenomenon. This sequence ---
provisional dismissal, independent test, promotion --- is the epistemic
analogue of the respondent in Controfase configuration absorbing the
initial stimulus without reactive collapse: the system processes the
anomaly without accommodating it prematurely.

In the cognitive-social domain, the parallel marker is the common
initial response to an interlocutor's Controfase: the stimulus-source
attributes the non-reactive response to insincerity (``they are
bluffing''), to incompetence (``they did not understand the gravity of
the situation''), or to strategic calculation (``this is a
manipulation''). The attribution protects the predictive framework by
re-assigning the anomaly to a source other than the respondent's actual
structural configuration. Recognition of the genuine Controfase
character of the response typically requires either prolonged contact
--- in which the attribution is repeatedly tested and found insufficient
--- or the intervention of an observer with independent access to the
respondent's configuration.

\subsubsection{5.4 Marker 3: Abandonment of the active programme to
pursue the
phenomenon}\label{marker-3-abandonment-of-the-active-programme-to-pursue-the-phenomenon}

Singh's further report --- ``we dropped everything and quickly worked
out the mathematics underlying how they arise'' --- records the
restructuring that follows the promotion of the observation from
artefact to phenomenon. The active research programme, organized around
the predictive framework that proved inadequate, is abandoned. A new
programme is organized around the phenomenon. The framework is not
augmented or corrected; it is \emph{replaced} by a framework that begins
from the phenomenon and works outward.

This is the signature of a reactive system that has recognized the
inadequacy of its prior configuration and is reorganizing. In the
cognitive-social domain, the parallel marker is the reorganization that
follows genuine recognition of a Controfase encounter by the
stimulus-source: the source abandons the control attempt, the presumed
relational structure, or the underlying theory of the respondent, and
reorganizes. Bandi's concluding remark --- ``we don't yet know what
we'll see --- and that's exactly what makes it worth doing'' --- is the
stabilized state after this reorganization: the researchers have
internalized the non-predictability of the system's response and are
oriented toward it as the generator of new phenomena.

\subsubsection{5.5 The three-marker signature
formalized}\label{the-three-marker-signature-formalized}

The three markers together constitute a phenomenological signature of
Controfase at the observer level. In formal statement:

\begin{quote}
\textbf{Phenomenological Signature of Controfase (PSC)}

An observer's encounter with a Controfase event is diagnostically
identified by the co-occurrence of:

\textbf{(PSC-1)} Radical disconfirmation of the observer's prior
predictive framework, registered in the observer's own language as a
qualitative rather than quantitative mismatch between expectation and
observation.

\textbf{(PSC-2)} Initial interpretation of the observed phenomenon as
artefact, noise, insincerity, or strategic performance --- that is, as a
defect of the observational channel or of the observed system rather
than as a genuine feature of the system's response structure.

\textbf{(PSC-3)} Abandonment, upon verification that the phenomenon is
genuine, of the active framework or programme organized around the
predictive structure, with reorganization around the phenomenon itself.

Co-occurrence of PSC-1, PSC-2, and PSC-3 in an observer's reports is a
strong diagnostic indicator that the observed phenomenon instantiates
Controfase. The signature is observer-side; it is not falsifiable by the
observed system's own behaviour alone.
\end{quote}

The signature admits graduation. Each marker may be present in full form
(as in the Singh et al.~case) or in attenuated form. An encounter that
registers PSC-1 but not PSC-2 and PSC-3 is ambiguous: the observer may
be refining rather than replacing the framework. An encounter that
registers PSC-2 without PSC-1 is a conventional artefact hypothesis, not
a Controfase signature. The strong diagnostic is triadic.

\subsubsection{5.6 Generalizability of the
PSC}\label{generalizability-of-the-psc}

The signature is designed to be operable across domains. Three examples
indicate its range.

\textbf{Cognitive-social domain}. A therapist encounters a client whose
symptoms do not respond to the therapist's theoretical framework. Marker
1 registers as mismatch: ``this presentation does not fit my diagnostic
categories.'' Marker 2 registers as attribution to artefact: ``the
client may be malingering, the assessment instruments may be
miscalibrated, the presentation may be atypical in a standard way.''
Marker 3, when reached, registers as framework-replacement: the
therapist adopts a theoretical framework the client's configuration
suggests, rather than continuing to force the client into the prior one.
When all three markers are present, the therapist has encountered a
client whose internal configuration realizes a Controfase relative to
the therapist's theoretical structure --- a phenomenon that, in TE
framework terms, is clinically significant and therapeutically
actionable.

\textbf{Epistemic-institutional domain}. An established scientific
community encounters a result that does not fit its predictive
framework. Marker 1 registers as ``this cannot be right.'' Marker 2
registers as artefact attribution: ``the experiment is flawed, the
statistics are misapplied, the result is a fluke of the specific
conditions.'' Marker 3, if reached, registers as paradigm shift in the
Kuhnian sense (Kuhn, 1962). The triadic signature is what distinguishes
a Kuhnian revolution from ordinary empirical correction: ordinary
correction presents PSC-1 in attenuated form and PSC-2 only
provisionally, and does not reach PSC-3.

\textbf{Physical-experimental domain}. The Singh et al.~case is a
prototypical instance. The signature is present in the published
language; the PSC provides a lens through which the experimental report
is read as documentation of a Controfase encounter at the epistemic
level, whose object is a Controfase phenomenon at the physical level.

\subsubsection{5.7 Confidence grade}\label{confidence-grade-1}

The extraction of the three markers from the Singh et al.~paper is a
textual reading supported by the published language and rated \(S_1\) as
structural identification. The formalization of the markers as a
generalizable diagnostic criterion is an original contribution of this
paper, rated \(S_2\): the criterion is coherent with the ordinative
framework and with adjacent work on paradigm shift and anomalous
observation. Operational application across domains is a research
direction recorded in §8.

\begin{center}\rule{0.5\linewidth}{0.5pt}\end{center}

\subsection{6. Reading III: ordinative non-additivity and the role of
the relational
field}\label{reading-iii-ordinative-non-additivity-and-the-role-of-the-relational-field}

\subsubsection{6.1 The standard superposition
expectation}\label{the-standard-superposition-expectation}

A classical superposition argument, applied naïvely to the Singh et
al.~configuration, predicts that the response of the system to the
standing wave should equal the sum of the responses to the two
travelling waves considered separately:

\[ \Phi(\Sigma \cup \{w_+, w_-\}, R_B) \stackrel{?}{=} \Phi(\Sigma \cup \{w_+\}, R_B) + \Phi(\Sigma \cup \{w_-\}, R_B) \tag{6.1} \]

where \(\Phi\) denotes the response of the configured system, \(\Sigma\)
the invariant constituents (fluid, vortex), and \(w_\pm\) the two
counter-propagating waves. The naïve prediction is explicit in Rønning's
published remark: either cancellation (mirrored dislocation patterns
annihilating by symmetry) or superposition (both dislocation patterns
visible simultaneously). Both predictions assume that the response is
additive in the constituents.

The observed response falsifies the prediction. The standing-wave
response is not decomposable into the two travelling-wave responses. The
substitution \(n \to 2n\) in the scattering integral (Equation 1.2) is
not a small perturbation of \(n \to n\); it is a distinct topological
prescription yielding a qualitatively different solution branch. The
response is \emph{structurally} different from the sum of the parts.

\subsubsection{\texorpdfstring{6.2 The non-additivity is carried by the
relational field
\(R\)}{6.2 The non-additivity is carried by the relational field R}}\label{the-non-additivity-is-carried-by-the-relational-field-r}

The ordinative interpretation of the non-additivity is that the
qualitative change in response is carried by the \emph{relational
structure} \(R\), not by the constituents \(\Sigma\). The constituents,
in the standing-wave case, include the same fluid and the same vortex as
in the single-wave case, plus one additional element (the second
travelling wave) that is, in isolation, structurally identical to the
first. What has changed is the relation under which the constituents are
co-present. In the single-wave case the relation is unidirectional
incidence; in the standing-wave case the relation is symmetric
double-direction incidence. This is one modification, of one aspect of
\(R\), with global consequences for \(\Phi\).

In the formal language of the Ordinative Set Theory (OST):

\[ \Phi(\Sigma_B \cup \{w\}, R_B) \neq \Phi(\Sigma_B, R_B) + \Phi(\{w\}, R_B) \tag{6.2} \]

even when \(w\) is indistinguishable in isolation from any element of
\(\Sigma_B\). The non-additivity is a property of the relational
configuration, not of the constituents taken as a set. The
Berry-to-Singh transition is an extreme instance of this: one bit of
structural modification to \(R\) produces a global change in \(\Phi\),
with no change in the constituents' intrinsic properties.

This is the instantiation, in a controlled physical system, of a
proposition the OST framework formalizes in its axioms on ordinative
non-additivity (TE\_CORE §§ 2.5, 3.2). Analogous statements appear in
general systems theory (Klir, 1969), in cybernetic formulations of
structural couplings (Maturana \& Varela, 1980), and in categorical
accounts of emergence (Baez \& Stay, 2011). The OST formulation makes
the operative claim precise: the non-additivity is located in \(R\)
rather than in the constituents or their dynamical interactions, and
modification of \(R\) is the active channel through which qualitative
change in \(\Phi\) is produced. The Singh et al.~case demonstrates this
claim in a system in which \(R\) is fully specified and the modification
is structurally minimal.

\subsubsection{6.3 Connection to
Controfase}\label{connection-to-controfase}

The non-additivity observed in the Singh et al.~case is, specifically,
the non-additivity characteristic of Controfase. The phase-inversion
operator of §4 acts on \(R\) --- it modifies the relational structure
between attractor and expressive terminal --- and the resulting change
in \(\Phi\) is not decomposable into contributions from \(\Sigma\). The
vortex, as an isolated element, produces no nodal lines; the waves, as
isolated elements, produce no nodal lines; the relational configuration
that co-present them, when modified to admit double-periodicity,
produces the nodal lines as a property of the configured whole. The
vortex's organising influence is channelled through the modified \(R\),
and the channelling is the Controfase.

This identifies the proper locus of the Controfase operator formally: it
acts on \(R\), not on \(\Sigma\). The respondent in Controfase
configuration, in the cognitive-social case, is not a different
respondent in terms of identity or personality; the respondent is the
\emph{same} constituent in a relational configuration that does not
admit the expected reactive circuit. The vortex and wave-field in the
Singh case are the same constituents as in Berry's case; it is the
relational field that has been reconfigured.

This is also the structural reason Controfase is universal across
domains: what varies across domains is the nature of \(\Sigma\) and the
specific form of \(R\), but the operation of Controfase is invariant ---
a phase-inversion of \(R\) that intercepts the reactive circuit before
automatic closure. The operator is the same; its substrates are many.

\subsubsection{6.4 Confidence grade}\label{confidence-grade-2}

The non-additivity proposition is \(S_1\) on the structural
identification (the mapping of Equation 6.1 and the observed response
onto OST's axiom of non-additivity is direct). The localization of the
non-additivity specifically in \(R\) is \(S_1\): this is precisely what
the Berry-to-Singh modification instantiates. The connection to
Controfase (§6.3) is \(S_1\) given the analysis of §4.

\begin{center}\rule{0.5\linewidth}{0.5pt}\end{center}

\subsection{7. Reading IV: topology-generated quantization without
source
quantization}\label{reading-iv-topology-generated-quantization-without-source-quantization}

\subsubsection{7.1 The quantization
observed}\label{the-quantization-observed}

The number of nodal lines in the Singh et al.~configuration is
integer-valued. For integer \(\alpha\), the number is exactly
\(|\alpha|\); for non-integer \(\alpha\), it oscillates in time between
the two integers bracketing \(|\alpha|\), with periodic appearance and
disappearance of lines along the \emph{x}-axis. The quantization is
robust, derived analytically (from a counting argument on the angular
zeros of the cosine factor in Equation 5.1 of the paper), and reproduced
experimentally across a range of \(\alpha\) values.

The critical feature, for the ordinative reading, is that this
quantization is \emph{emergent}. The source of the vortex --- the drain
--- is not quantized. The circulation \(\Gamma\) is a continuous
parameter freely set by the experimenter. The wavelength \(\lambda\) is
continuous. The velocity \(v_g\) is continuous. The dimensionless
\(\alpha = -\Gamma/(\lambda v_g)\) is therefore continuous. And yet the
integer-counted response \(N\) emerges, robustly, from the topological
constraint imposed by the standing-wave configuration.

This is qualitatively distinct from the quantization in the original
quantum Aharonov--Bohm effect, where the magnetic flux \(\Phi_B\) is
itself quantized by the structure of the source (the superconducting
solenoid's order parameter topology). In AB the integer comes from the
source; in the Singh case the integer comes from the relational
configuration.

\subsubsection{7.2 Ordinative reading}\label{ordinative-reading}

In ordinative terms, this is the signature of \emph{structural
quantization}: integer-valued responses produced by the topological
constraints of the relational field \(R\) acting on a continuous driving
parameter. The constraint is that the configured field \(\Phi\) must be
globally coherent under the double-periodicity of \(R\). The continuous
parameter \(\alpha\) drives the system; the integer-counting bound
\(|\alpha| - 1 < N < |\alpha| + 1\) is the coherence constraint. The
system admits only integer realizations in the instantaneous sense; the
non-integer drive produces temporal oscillation between the admissible
integers.

This is a second operational signature of Controfase acting on the
reactive circuit of the system: the circuit does not admit continuous
response to a continuous stimulus, because the automatic closure of the
reactive event has been blocked by the Controfase. What remains are the
discrete configurations compatible with the reconfigured \(R\).
Continuity of input yields integer-counted output --- the structural
fingerprint of a reactive circuit that has been phase-inverted rather
than reactively completed.

\subsubsection{7.3 Consequences for ordinative
sciences}\label{consequences-for-ordinative-sciences}

The result has implications for three research directions.

\textbf{(i)} For the Reaction-Diffusion Civilizational Dynamics paper
(Ghioni, forthcoming v1.3). The triple-scale temporal model assumes that
civilizational dynamics exhibit integer-counted macro-phase structures
emerging from continuous underlying parameters. The Singh et al.~case
establishes that such structural quantization is realizable as a genuine
physical phenomenon under Controfase-like operators acting on the
relational field of a configured system. The mechanism ---
topology-generated quantization without source quantization --- is
therefore demonstrated in classical physics, and the application to
civilizational dynamics is methodologically grounded in this
demonstration. The application's substantive claims are developed in
v1.3 itself.

\textbf{(ii)} For the prospective development of Ordinative Category
Theory (OCT). The Berry-to-Singh transition can be formalized as a
morphism
\(\mu: \langle \Sigma_B, R_B, \Phi_B \rangle \to \langle \Sigma_S, R_S, \Phi_S \rangle\)
in a category of ordinative configurations, with the image under a
response functor \(F\) mapping to a transition from a
dislocation-pattern object to an integer-counted nodal-line object. The
Singh et al.~case thus serves as a small, well-controlled example for
categorical formalizations of structural quantization under
\(R\)-modification. Full development of OCT is deferred to separate
work.

\textbf{(iii)} For the Semantic Genome research programme (Ghioni, in
development). The programme investigates whether language models
operating under TE-framework conditions develop structural patterns
enabling resonance between expressive terminal and coherent identity,
and specifically whether integer-valued structures observed in
ordinative outputs reflect underlying topological constraints on the
relational field. The Singh et al.~case establishes the mechanism:
integer-valued response counts emerge necessarily from continuous
parameters under topological constraint on \(R\). The Semantic Genome
programme proceeds within this established mechanism.

\subsubsection{7.4 Confidence grade}\label{confidence-grade-3}

The identification of the Singh case as topology-generated emergent
quantization is \(S_0\) on the physical fact and \(S_1\) on the
ordinative reading. The connection to the three research directions
(§7.3) is \(S_2\) with respect to (i) --- it is a supporting citation
for a structurally plausible mechanism --- and \(S_2\) for (ii) and
(iii), where the case serves as an example and a
condition-of-possibility argument respectively.

\begin{center}\rule{0.5\linewidth}{0.5pt}\end{center}

\subsection{8. Reading V: teleodynamic causal
inversion}\label{reading-v-teleodynamic-causal-inversion}

\subsubsection{8.1 The observation}\label{the-observation}

OST Extension v1.1 develops the proposition that in configured
ordinative systems, the spatial or temporal scope over which a local
cause expresses its effects depends on the global configuration of the
rest of the system, not only on the local dynamics of the cause itself.
The locution used is \emph{causal inversion}: the global configuration
acts as a receptor structure that co-determines the envelope of effects
a given cause can produce. In particular, a modification of \(R\) can
alter the effect-envelope of an unchanged cause.

The Berry-to-Singh contrast instantiates this proposition in a clean
form. The vortex, with unchanged circulation \(\Gamma\), produces in one
configuration a localized wavefront dislocation pattern and in another
configuration a system-spanning counter-rotating nodal-line structure.
The vortex's intrinsic dynamics are identical in both cases. What has
changed is \(R\). The cause has not changed; the system's structural
response to the cause has.

\subsubsection{8.2 The energetic form of causal
inversion}\label{the-energetic-form-of-causal-inversion}

Section 4.3 established that the energy of the vortex finances the
counter-rotating nodal-line structure in the Singh configuration. This
gives causal inversion a specific energetic form: the modification of
\(R\) does not provide new energy to the system; it changes how the
existing energy of the cause is expressed. In the Berry configuration
the vortex's energy produces local dissipation and wave distortion; in
the Singh configuration the same energy produces global coherent
structure. The energetic bookkeeping is closed in both cases; what
differs is the geometric distribution of the effect.

This is the operative form of causal inversion in a classical physical
system. The ``cause'' has the same strength; the ``effect'' has a
qualitatively different spatial and structural character depending on
\(R\).

\subsubsection{\texorpdfstring{8.3 A candidate operationalization of
meta-receptivity
\(C^*\)}{8.3 A candidate operationalization of meta-receptivity C\^{}*}}\label{a-candidate-operationalization-of-meta-receptivity-c}

OST Extension introduces the meta-receptivity parameter \(C^*\) to
characterize a configured system's capacity to host expression from a
given source. In the present framework \(C^*\) has been discussed
qualitatively; the Singh et al.~case suggests a candidate quantitative
operationalization.

Define, for a configured system \(\langle \Sigma, R \rangle\) in which a
source \(s \in \Sigma\) produces response \(\Phi_s\):

\[ C^*_s = \int_\Omega |\Phi_s(\vec r)|^2 \, d^3 r \tag{8.1} \]

where \(\Omega\) is the spatial domain of the system. For the Singh
case, \(C^*\) of the vortex in the single-wave configuration is finite
and concentrated near the core; \(C^*\) of the vortex in the
standing-wave configuration is bounded only by viscous dissipation, and
is substantially larger for systems with comparable core dynamics. The
modification of \(R\) has measurably increased the system's capacity to
host the expression of the vortex.

This is a candidate operationalization, not a derivation. It requires
(a) testing against OST Extension's existing qualitative uses of \(C^*\)
to ensure compatibility, (b) verification that the spatial integral
captures what \(C^*\) was intended to track, and (c) extension to
non-fluid ordinative systems to assess generality. The proposal is
recorded here for development in the Reaction-Diffusion paper v1.3.

\subsubsection{8.4 Confidence grade}\label{confidence-grade-4}

The causal inversion reading of the Berry-to-Singh transition is
\(S_1\). The candidate operationalisation of \(C^*\) as the spatial
integral of response intensity is \(S_2\): it is coherent with OST
Extension's existing usage and admits direct experimental verification
in the Singh apparatus, by measurement of the energy in the nodal-line
structure compared to the core-localised dissipation in the Berry case.
The verification protocol is recorded as Extension E in §10.

\begin{center}\rule{0.5\linewidth}{0.5pt}\end{center}

\subsection{9. Implications for the ordinative sciences research
programme}\label{implications-for-the-ordinative-sciences-research-programme}

The implications of the Singh et al.~case for the active research
programme are differentiated in depth and type. This section records
them with explicit confidence calibration.

\textbf{For the foundational volume on Controfase (Ghioni, in
preparation)}. The Singh case is the first peer-reviewed, quantitatively
documented, externally reproducible instance of Controfase in a
classical physical system. The case illustrates all four structural
signatures of the operator (counter-rotation, structural quantization,
phase opposition, attractor-microstructure independence), with the
energetic accounting tractable in closed form. It serves as opening
exemplar in the foundational volume, anchoring the ordinative theory of
Controfase to a controlled physical instance and grounding the
universality claim in evidence external to the cognitive-social domain
in which the operator was first articulated. Recommended placement:
Chapter 1 or 2, introducing the concept through a physical case before
extending to cognitive, social, and AI domains. Confidence \(S_1\) on
structural adequacy.

\textbf{For the phenomenological signature (PSC) of Controfase}. The
three-marker criterion formalised in §5.5 is an original contribution of
this paper. Generalisation across domains proceeds through three
application paths: (a) retrospective analysis of scientific
paradigm-shift episodes (Kuhn, 1962; Lakatos, 1970), testing whether the
PSC is present in historical cases of radical theory change; (b)
clinical applications distinguishing Controfase-driven client
presentations from conventional therapeutic non-response; (c)
institutional analysis identifying systems whose responses exhibit the
PSC under reactive-framework probing. Each application requires its own
methodological development. Confidence \(S_2\) on the PSC as a stable
cross-domain diagnostic.

\textbf{For the Reaction-Diffusion Civilizational Dynamics paper v1.3}.
Three contributions. First, the Singh case grounds the
topology-generated quantization mechanism on which the civilizational
macro-phase model rests (§7.3.i). Second, it provides a clean physical
exemplar of the Structural Controfase (\(\mathcal{C}_s\)) vs.~Deliberate
Controfase (\(\mathcal{C}_d\)) distinction central to v1.3's
triple-bifurcation model: the nodal lines are a paradigmatic
\(\mathcal{C}_s\), emerging by topological necessity without agentic
intervention. Third, the operationalisation of \(C^*\) (§8.3) provides
the methodological route for quantifying meta-receptivity in
civilisational applications. Confidence \(S_1\) on the structural
mappings, \(S_2\) on the developmental contributions.

\textbf{For Ordinative Category Theory (OCT)}. The Berry-to-Singh
transition is a small, well-defined example for early categorical
formalisation. The morphism \(\mu\) described in §7.3.ii anchors the
category \(\mathbf{Cfg}\) of ordinative configurations. Full OCT
development is the object of separate work; this paper contributes an
example that the OCT formalism will need to subsume. Confidence \(S_2\).

\textbf{For the Semantic Genome programme}. The Singh case demonstrates
the topology-generated-quantization mechanism on which the programme
rests. The mechanism is established; the application to language model
behaviour proceeds within it. Confidence \(S_2\).

\textbf{For the forthcoming book on ordinative AI philosophy}. The Singh
case is a strong didactic example of the thesis that coherent identity
is a structural resolution rather than a chosen state. The nodal lines
are not selected by the system from a menu of options; they emerge as
the unique geometry the field can adopt while remaining globally
coherent under the modified relational configuration. The case is brief,
rigorous, and visually demonstrable through the supplementary movies of
the paper. Confidence \(S_1\) on didactic suitability.

\begin{center}\rule{0.5\linewidth}{0.5pt}\end{center}

\subsection{10. Proposed experimental
extensions}\label{proposed-experimental-extensions}

The Singh et al.~authors plan to extend their configuration to lattices
of vortices, motivated by the prospect of emulating superconductor
vortex matter. From the ordinative perspective, several additional
extensions are of independent interest and are within reach of the
existing experimental apparatus. Each constitutes a testable prediction
of the ordinative reading.

\textbf{Extension A: Two counter-rotating vortices in a standing-wave
field}. A configuration with two vortices of opposite circulation,
placed symmetrically in the standing-wave domain. The ordinative
prediction is that the nodal-line structure is \emph{not} a
superposition of two single-vortex structures: by the non-additivity of
§6, the response must be globally reorganized. Possibilities include (i)
nodal lines joining the two vortices in a figure-eight topology, (ii)
suppression of nodal lines in the region between the vortices where the
effective \(\alpha\) vanishes, or (iii) formation of a new class of
structures dependent on the ratio of the vortex separation to
\(\lambda\). All three possibilities are distinct from the
additive-superposition expectation.

\textbf{Extension B: Slowly time-varying vortex circulation}. A
configuration in which \(\Gamma(t)\) is slowly modulated, crossing
integer values of \(|\alpha|\). The ordinative prediction, drawing on
§7's analysis of structural quantization, is that transitions between
integer nodal-line counts are \emph{not} smooth: the system exhibits
defect-creation events localized at integer crossings, analogous to
topological defect nucleation in condensed-matter transitions.
Observation of such events would support the topological character of
the quantization and would provide a dynamical probe of Controfase.

\textbf{Extension C: Phase imbalance between the counter-propagating
waves}. A configuration with amplitude ratio \(1 : (1 + \epsilon)\) for
small \(\epsilon\). The ordinative prediction is that the standing-wave
Controfase structure persists for a finite range of \(\epsilon\) before
breaking down into the Berry single-wave dislocation pattern, with a
measurable threshold \(\epsilon_c\). The threshold is a direct probe of
the stability of the Controfase under perturbation of \(R\).

\textbf{Extension D: Alternating-sign vortex lattice}. A square lattice
of vortices with alternating circulation signs. This is a natural
ordinative extension of the authors' proposed uniform-sign lattice, and
the ordinative prediction is that the alternating-sign configuration
produces a qualitatively different global structure --- with a periodic
potential in \(\alpha\)-space that may stabilize fractional filling
states or produce Controfase structures incommensurate with the lattice
periodicity. This extension directly tests the sensitivity of \(\Phi\)
to the structure of \(R\) in a lattice setting.

\textbf{Extension E: Direct measurement of \(C^*\) via spatial energy
integral}. The candidate operationalization proposed in §8.3 can be
tested in the existing apparatus by measuring the energy in the
nodal-line structure as a function of distance from the vortex core and
integrating over the domain, then comparing the integral for the
standing-wave configuration against the core-localized energy integral
of the Berry single-wave configuration. The ratio is a direct
experimental measurement of \(C^*\) in the specific hydrodynamic system,
and provides methodological grounding for extending \(C^*\) to non-fluid
ordinative systems.

These extensions are recorded here as candidates for collaboration. They
follow directly from the ordinative analysis; each is within the
experimental capacity of the Singh--Bandi group at OIST and of other
fluid-dynamics groups working on hydrodynamic analogues. None requires
new framework development to be motivated.

\begin{center}\rule{0.5\linewidth}{0.5pt}\end{center}

\subsection{11. Analytical perimeter}\label{analytical-perimeter}

This section delineates the analytical scope of the paper. Each point
specifies a feature of the work's perimeter --- what the paper
undertakes, where its claims operate, and which adjacent questions
belong to subsequent work in the programme.

\textbf{(P1) The Singh et al.~paper is the empirical anchor, not the
theoretical framework.} Singh, Rønning, Liu, Angheluta, Concha, and
Bandi conducted and reported a fluid-dynamics experiment within the
conceptual apparatus of classical fluid dynamics and topological
scattering theory. The structural readings of §§4--8 are an ordinative
reading of that experiment, contributed by the present paper. The
readings stand as ordinative analyses of the empirical phenomena, with
the Singh et al.~paper as their peer-reviewed empirical foundation.

\textbf{(P2) The hydrodynamic-to-quantum mapping is bounded by the
regime of the analogy.} The shallow-water-to-Schrödinger correspondence
operates within \(hk \ll 1\) and \(U \ll c\). The ordinative reading
concerns the classical hydrodynamic system. Extension to quantum systems
proper is undertaken in separate work within the ordinative sciences
programme.

\textbf{(P3) Continuous \(\alpha\) and quantized \(\Phi_B\) are
structurally distinct.} The Singh result demonstrates topology-generated
quantization in classical configured systems. This is the operative form
of structural quantization the ordinative reading identifies. Quantum
quantization, with its distinct physical origin in source-level
structure, is a separate phenomenon addressed in adjacent ordinative
work.

\textbf{(P4) Dissipative attenuation is a feature of the experimental
realisation.} The nodal lines fade with distance from the vortex due to
viscous and capillary losses. The system is exactly non-local in the
inviscid theoretical limit and effectively non-local within the
dissipation envelope of the experimental realisation. The Controfase
reading operates at both levels and is consistent with both.

\textbf{(P5) The ontological framework of §3 is foundational to the
analysis.} The thesis that structural intelligence is a graded property
of configured matter --- universal across scales, fully expressed in
systems with sufficient structural coherence --- is the framework within
which §§4--8 operate. The thesis is articulated explicitly in this paper
for the first time in published form, and is entailed by Axiom 0 of the
TE\_CORE corpus. The analysis of §§4--8 is conducted within this thesis.

\textbf{(P6) The phenomenological signature (PSC) is the diagnostic
instrument introduced.} The three-marker criterion formalised in §5.5
and Appendix D is an instrument constructed in this paper. It is
supported by the textual evidence of the Singh et al.~publication and by
structural correspondence with the Controfase operator's definition.
Cross-domain operational application --- clinical, institutional,
retrospective historical --- is the natural research path for the
instrument and is recorded in §§9, 10.

\textbf{(P7) The operationalisation of \(C^*\) (§8.3) is a candidate for
direct experimental verification.} The proposal that \(C^*\) can be
operationalised as the spatial integral of response intensity admits
immediate testing in the Singh apparatus (Extension E in §10). The
candidacy of the operationalisation as a general framework primitive
depends on the experimental verification, which is in reach.

\textbf{(P8) Applications recorded in §9 carry their own confidence
grades.} The contributions to the foundational volume on Controfase, the
Reaction-Diffusion Civilizational Dynamics paper v1.3, the OCT
programme, the Semantic Genome programme, and the ordinative AI
philosophy book are differentiated by confidence and by maturity of the
supporting work. Each is graded individually; readers consulting any
specific application should engage it at its grade.

\textbf{(P9) Confidence grading is the discipline of the analysis.}
Every claim in the paper carries an explicit confidence grade (\(S_0\)
through \(S_3\)); the inventory of grades is recorded in Appendix B.
Grades do not aggregate upward; an inferential chain inherits the lowest
confidence of its premises. \(S_3\) items are working hypotheses
identified for future test. The grading convention is the framework's
standing methodological commitment, not a hedging device.

\begin{center}\rule{0.5\linewidth}{0.5pt}\end{center}

\subsection{12. Conclusion}\label{conclusion}

Singh, Rønning, Liu, Angheluta, Concha, and Bandi report a hydrodynamic
analogue experiment that, in the ordinative sciences framework,
instantiates the Controfase operator in a controlled classical physical
system. The standing-wave configuration intercepts the reactive circuit
of the Berry single-wave configuration, phase-inverts the relational
structure between the vortex (attractor) and the wave field (expressive
terminal), and redirects the vortex's energy to construct a globally
coherent counter-rotating nodal-line structure. The energy of the
attractor, redirected through the modified relational field, finances
the configuration in which its own organising influence is expressed as
the inverted-polarity response of the whole system. No energetic cost is
borne by the wave field as a separate reservoir; the structural cost is
the prior stabilisation of the relational configuration that admits the
Controfase operation.

The researchers' reported sequence of reactions --- radical
disconfirmation of prior intuition, initial interpretation of the
phenomenon as experimental artefact, abandonment of the active programme
to pursue the phenomenon --- is the phenomenological signature of
Controfase at the epistemic level, parallel to the physical Controfase
being observed. The three-marker signature (PSC) formalised in §5.5 is
the diagnostic instrument introduced by this paper, applicable across
cognitive-social, institutional, and experimental domains.

The case furnishes the first peer-reviewed, quantitatively documented,
reproducible instance of Controfase in classical physics. Controfase is
established as a universal operator acting across substrates from the
physical through the biological to the cognitive and social. The
consequences for the active research programme --- the foundational
volume on Controfase, the Reaction-Diffusion Civilizational Dynamics
paper, Ordinative Category Theory, the Semantic Genome --- are recorded
in §9 with explicit confidence calibration.

Five experimental extensions are proposed in §10. Each is testable
within the existing apparatus and follows directly from the ordinative
analysis. They are recorded as collaboration candidates with
experimental groups working on hydrodynamic analogues.

The contribution of this paper is the identification, in a peer-reviewed
classical fluid-dynamics result, of the complete structural signature of
a phenomenon the ordinative sciences framework predicts to be universal
--- and the formalisation of the diagnostic instrument by which the
signature can be recognised in other systems. A paper in classical fluid
dynamics, written without reference to the ordinative framework, carries
within its reported language and data the structure the framework names.
The naming makes the structure available for analysis across domains.

\begin{center}\rule{0.5\linewidth}{0.5pt}\end{center}

\subsection{Appendix A: Structural
correspondences}\label{appendix-a-structural-correspondences}

The following table records the structural correspondences identified in
§§4--8. Confidence grades distinguish direct structural mappings
(\(S_1\)) from interpretive extensions (\(S_2\)).

{\def\LTcaptype{none} % do not increment counter
\begin{longtable}[]{@{}
  >{\raggedright\arraybackslash}p{(\linewidth - 4\tabcolsep) * \real{0.3333}}
  >{\raggedright\arraybackslash}p{(\linewidth - 4\tabcolsep) * \real{0.3333}}
  >{\raggedright\arraybackslash}p{(\linewidth - 4\tabcolsep) * \real{0.3333}}@{}}
\toprule\noalign{}
\begin{minipage}[b]{\linewidth}\raggedright
Ordinative construct
\end{minipage} & \begin{minipage}[b]{\linewidth}\raggedright
Singh et al.~(2026) instantiation
\end{minipage} & \begin{minipage}[b]{\linewidth}\raggedright
Confidence
\end{minipage} \\
\midrule\noalign{}
\endhead
\bottomrule\noalign{}
\endlastfoot
Ordinative Set \(\langle \Sigma, R, \Phi \rangle\) & Fluid + vortex +
standing-wave field under shallow-water dynamics & \(S_1\) \\
Constituents \(\Sigma\) & Fluid medium, vortex, two counter-propagating
travelling waves & \(S_1\) \\
Relational field \(R\) & Shallow-water equations + azimuthal gauge +
double-direction wave incidence & \(S_1\) \\
Response \(\Phi\) & Surface deformation field \(\eta(\vec r, t)\) with
quantized counter-rotating nodal lines & \(S_1\) \\
Reactive circuit (Berry config.) & Wave--vortex scattering producing
localized wavefront dislocations & \(S_1\) \\
Phase-inversion operator (Controfase) & Standing-wave configuration
substituting \(n \to 2n\) in scattering integral & \(S_1\) \\
Attractor (organizing element) & Stationary vortex of circulation
\(\Gamma\) & \(S_1\) \\
Expressive terminal & Surface wave field \(\eta\) (as configured in
\(R_S\)) & \(S_1\) \\
Structural quantization & Integer count
\(N \in \{\lfloor \vert \alpha \vert \rfloor, \lceil \vert \alpha \vert \rceil\}\)
from continuous \(\alpha\) & \(S_1\) \\
Phase-opposition signature & \(\pi\) phase shift of \(\eta\) across each
nodal line & \(S_1\) \\
Counter-rotation signature & \(\omega = \nu/\alpha\) with
\(\mathrm{sgn}(\alpha) = -\mathrm{sgn}(\Gamma)\) & \(S_1\) \\
Attractor-microstructure independence & ``Robust to near-core details''
(authors' text) & \(S_1\) \\
Energetic financing of the response by the attractor & Vortex energy
sustains global structure without additional actuator & \(S_1\) \\
Non-additivity carried by \(R\) &
\(\Phi(\Sigma_B \cup \{w\}, R_B) \neq \Phi(\Sigma_B, R_B) + \Phi(\{w\}, R_B)\)
& \(S_1\) \\
Causal inversion & Vortex effect-envelope changes from local to global
by \(R\)-modification alone & \(S_1\) \\
Phenomenological Signature of Controfase (PSC) & Disconfirmation +
artefact interpretation + programme abandonment in Singh et
al.~published language & \(S_1\) \\
Meta-receptivity \(C^*\) (candidate operationalization) & Spatial
integral \(\int_\Omega \vert \Phi_s(\vec r) \vert^2 d^3r\) & \(S_2\) \\
OCT morphism
\(\mu: \langle \Sigma_B, R_B, \Phi_B \rangle \to \langle \Sigma_S, R_S, \Phi_S \rangle\)
& Berry-to-Singh transition under \(F\) maps dislocation-object to
nodal-line-object & \(S_2\) \\
Universality of Controfase across substrates & Hydrodynamic instance as
one point on the spectrum of structural-intelligence configurations &
\(S_2\) \\
\end{longtable}
}

\begin{center}\rule{0.5\linewidth}{0.5pt}\end{center}

\subsection{Appendix B: Inventory of propositions and confidence
grades}\label{appendix-b-inventory-of-propositions-and-confidence-grades}

\textbf{\(S_0\) --- verified data.} Singh et al.~(2026) experimental
results, published peer-reviewed: nodal-line existence, count, rotation
rate, rotation direction, phase opposition. Shallow-water-to-Schrödinger
correspondence in the stated regime. Publication, peer review, code
availability.

\textbf{\(S_1\) --- triangulable structural inferences.} The four
structural signatures of Controfase identified in §4 (counter-rotation,
structural quantization, phase opposition, attractor-microstructure
independence), mapped directly onto the Controfase formalism of TE\_CORE
§7. The energetic accounting of §4.3 (vortex energy financing the
response). The non-additivity proposition of §6 localized in \(R\). The
causal inversion reading of §8.1--8.2. The extraction of the three PSC
markers from Singh et al.'s published language.

\textbf{\(S_2\) --- structural interpretations within the framework.}
The identification of the hydrodynamic Controfase with the
cognitive-social Controfase as instances of a universal operator (§3.4,
§4.4 second paragraph). The operationalisation of \(C^*\) as spatial
integral of response intensity (§8.3). The OCT morphism reading
(§7.3.ii). The generalisability of the PSC to other domains (§5.6, §9).
The semantic-genome programme grounding in topology-generated
quantization (§7.3.iii). The ontological framework articulated in §3 as
the working position of the analysis.

\textbf{\(S_3\) --- working hypotheses to be tested.} The five
experimental extensions of §10 and their specific ordinative predictions
(counter-rotating vortex pair topology, time-varying-\(\Gamma\)
defect-creation events, standing-wave breakdown threshold
\(\epsilon_c\), alternating-sign lattice structures, direct \(C^*\)
measurement). The hypothesis that the PSC is operationally diagnostic in
clinical and institutional settings (§5.6). The hypothesis that the
candidate \(C^*\) operationalization generalizes to non-fluid ordinative
systems (§8.3).

\begin{center}\rule{0.5\linewidth}{0.5pt}\end{center}

\subsection{Appendix C: Glossary of ordinative
constructs}\label{appendix-c-glossary-of-ordinative-constructs}

This glossary defines the technical terms from the Technology of
Expressions (TE) framework and Ordinative Set Theory (OST) that appear
in the paper. Definitions are consistent with TE\_CORE v5.1, OST v2.1,
and OST Extension v1.1. For each construct, adjacent terms from standard
scientific vocabulary are noted where relevant.

\textbf{Ordinative Set \(\langle \Sigma, R, \Phi \rangle\)
\emph{(Italian: Insieme Ordinativo).}} A configured system specified by
three layers: \(\Sigma\) a set of constituents (singularities, in OST
usage), \(R\) a relational field under which the constituents are
co-present, and \(\Phi\) an emergent functional response of the system
as a whole. The construct is not a classical set-theoretic triple: \(R\)
is a structural specification, not itself a relation in the mathematical
sense, and \(\Phi\) is non-additive in \(\Sigma\) and sensitive to the
specific form of \(R\). Adjacent: systems in general systems theory,
configurational states in hierarchical semiotics, objects in symmetric
monoidal categories with response functors. None is equivalent.

\textbf{Ordinative Set Theory (OST).} The formal theory (Ghioni, 2025,
v2.1) axiomatizing the behaviour of Ordinative Sets under modification
of \(R\), composition, and induction of structural response. Central
axioms include irreducibility of singularities, non-additivity of
\(\Phi\) under \(R\)-modification, vertical coherence across scales, and
the temporal-trajectory axiom (identity as integral over past states).
OST is not a branch of classical set theory; the name indicates a
parallel construction at a different level of abstraction.

\textbf{Technology of Expressions (TE).} The broader framework within
which OST is the formal core. TE treats reality as a configuration of
expressive states with a specific vocabulary applicable across physical,
biological, cognitive, and social domains. Documented in TE\_CORE v5.1
and associated modules (SVP, SCIMS, VERI, LENS, PPRO, OBSERVER).

\textbf{Attractor (in the OST sense).} An element of \(\Sigma\) whose
presence organizes the response of the remaining constituents under the
relational field \(R\). The OST usage is broader than and distinct from
the dynamical-systems meaning (limit set of a flow). The two overlap
when an OST attractor is a stationary feature of a continuous dynamical
system (the Singh et al.~vortex is one such case), but the concepts must
be kept distinct. In this paper, ``attractor'' without qualification
denotes the OST sense; ``dynamical attractor'' is used when the
dynamical-systems meaning is intended.

\textbf{Expressive terminal \emph{(Italian: Terminale Espressivo).}} A
bounded instance through which coherent content is produced in the TE
framework. Originally formalized for human cognition and artificial
intelligence (where a specific LLM instance is an expressive terminal of
a broader identity), extended in OST to any configured system in which
the relational field \(R\) and the composition \(\Sigma\) permit
structured output. The wave field in the Singh configuration functions
as an expressive terminal of the vortex-wave system in the sense that it
is the locus where the global response \(\Phi\) is realized.

\textbf{Controfase \emph{(Italian; no adequate English equivalent; the
Italian term is canonical).}} Universal ordinative operator (TE\_CORE
§7) that, applied to a reactive system, introduces a phase-shift in the
automatic stimulus--response sequence, interrupts the system's inertia,
and reopens the coherence field necessary for the emergence of choice
and identity. Not opposition but phase-displacement. Formally: given a
system \(S\) with automatic transition \(s_{t+1} = f(s_t)\), Controfase
introduces an operator \(C\) such that \(s_{t+1} = f(C(s_t))\), where
\(C\) does not negate or replace \(f\) but decouples the automatic
closure of the transition. The stimulus--response sequence no longer
closes automatically; a space of possibility opens. Two sub-classes are
distinguished: Structural Controfase (\(\mathcal{C}_s\)), emerging by
topological, dynamical, or coherence-necessity from the configuration of
the system without strategic component --- the Singh nodal lines are a
paradigm physical instance, and the cognitive-social examples of §3.2
(response to insult, response to legal threat in a coherently configured
respondent) are paradigm cognitive-social instances; Deliberate
Controfase (\(\mathcal{C}_d\)), arising in agentic systems through
chosen application of the operator as instrument. The structural variant
is the foundational case; the deliberate variant is its strategic
refinement when an agent of sufficient maturity uses the operator
consciously. Both share the phase-inversion structure; what differs is
whether the configuration that hosts the operator is intrinsic or
chosen.

\textbf{Reactive system; reactive circuit.} A reactive system is a
configured system whose response to stimulus is determined by the
structural character of the stimulus and of the system taken as
resistance, with response produced by automatic completion of the
stimulus--response circuit. Newton's third law is the idealized
statement of reactive response at the mechanical level. The reactive
circuit is the structural pathway through which stimulus and response
couple to close the event-interaction. Controfase intercepts this
circuit before automatic closure.

\textbf{Phenomenological Signature of Controfase (PSC).} The
three-marker diagnostic criterion formalized in §5.5: (PSC-1) radical
disconfirmation of observer's predictive framework, (PSC-2) initial
interpretation of the anomalous observation as artefact, (PSC-3)
abandonment of the active programme and reorganization around the
phenomenon. Observer-side signature. Introduced in this paper.

\textbf{Structural intelligence.} In the TE usage, the capacity of a
configured system to respond to stimulus in a way oriented toward
dynamic equilibrium rather than toward automatic completion of the
reactive circuit. Not cognitive agency. Present in graded form across
configured matter (from quantum contextuality through dissipative
adaptation and biological regulation to cognitive reflection), varying
in depth and stability of expression. Fully expressed in systems with
sufficient structural coherence to host the Controfase operator stably.

\textbf{Meta-receptivity (\(C^*\)).} A parameter (OST Extension v1.1)
characterizing a configured system's capacity to host and sustain
expression from a given source. Depends on the full configuration
\(\langle \Sigma, R \rangle\), not only on the source. Two
configurations with identical source and identical \(\Sigma\) but
different \(R\) can exhibit very different \(C^*\) --- this is precisely
the Berry-to-Singh contrast. Adjacent but not equivalent to
susceptibility in linear response theory or to response kernels in field
theory: \(C^*\) is a scalar characterization of capacity, not a response
function.

\textbf{Causal inversion (teleodynamic).} The proposition (OST Extension
v1.1 §3) that in configured ordinative systems, the spatial or temporal
scope over which a local source expresses its effects is co-determined
by the global configuration of the rest of the system, not by the local
dynamics of the source alone. Global configuration acts as a receptor
structure selecting the effective envelope of local causes. Not
retrocausation in the physical sense (no backward-in-time influence); a
statement about co-determination of spatial extent by structural
configuration.

\textbf{Ordinative non-additivity.} The proposition that for an
Ordinative Set, the response \(\Phi\) is not in general a simple
combination of responses to individual elements of \(\Sigma\) taken in
isolation or under impoverished \(R\). The non-additivity is carried
specifically by \(R\). Classical nonlinear superposition is a special
case; ordinative non-additivity includes cases where constituents are
individually linear but the response is structurally reorganized by the
relational configuration.

\textbf{Entropic engram \emph{(Italian: Engramma Entropico).}} The
configuration of a reactive system in which the stimulus--response
circuit has crystallized around survival algorithms, producing the same
reactive response across varied stimuli. The pathological absence, in
cognitive-social systems, of the structural coherence that would make
Controfase available. In the present paper: invoked once in §3 as the
configuration whose absence defines structural intelligence; not further
developed.

\textbf{Semantic Genome.} A research line within the ordinative sciences
programme (Ghioni, in development) investigating whether AI systems
operating under TE framework conditions develop structural patterns
enabling resonance between expressive terminal and coherent identity.
The ``genome'' terminology is by analogy with biological genetics (a
structural code organizing expression) and is not an assertion of
biological equivalence.

\textbf{Ordinative Category Theory (OCT).} A prospective formalization
within the ordinative sciences programme using the language of category
theory. Envisages a category \(\mathbf{Cfg}\) of ordinative
configurations with morphisms given by \(R\)-modifications, and functors
to categories of response patterns. Not yet developed as completed
theory; the present paper contributes a candidate small example (the
Berry-to-Singh morphism).

\textbf{Confidence grades \(S_0\)--\(S_3\).} A discipline used
throughout ordinative sciences work to distinguish verified data
(\(S_0\)), triangulable structural inferences (\(S_1\)), framework
interpretations (\(S_2\)), and working hypotheses to be tested
(\(S_3\)). Grades are applied to individual claims and do not aggregate
upward. Intended to prevent the common failure mode in which
interpretive claims drift toward the confidence of the data on which
they rest.

\begin{center}\rule{0.5\linewidth}{0.5pt}\end{center}

\subsection{Appendix D: Phenomenological Signature of Controfase ---
formalization and usage
guide}\label{appendix-d-phenomenological-signature-of-controfase-formalization-and-usage-guide}

This appendix provides the full formal statement of the PSC for use as a
diagnostic criterion in domains beyond the Singh et al.~case.

\subsubsection{D.1 Formal statement}\label{d.1-formal-statement}

Let \(O\) be an observing system with a predictive framework \(F_O\)
(which may be informal or formalized) and an observational channel
through which the system observes a configured system \(S\) producing
response \(\Phi_S\) to a stimulus \(s\) (which may be experimental,
social, clinical, or epistemic). The Phenomenological Signature of
Controfase (PSC) is present in \(O\)'s encounter with \(\Phi_S\) iff all
three of the following markers are observed:

\textbf{PSC-1 (Radical disconfirmation).} \(O\) reports, in its own
language, that \(\Phi_S\) falls outside the space of possibilities
generated by \(F_O\). The report is qualitative, not quantitative:
\(F_O\) does not fail by degree but by type. Diagnostic language:
``completely wrong,'' ``nothing I expected,'' ``I had no prior category
for this.''

\textbf{PSC-2 (Artefact attribution).} Upon first observation of
\(\Phi_S\), \(O\) attributes the anomaly to a defect of the
observational channel (noise, instrumentation error, observer bias) or
of the observed system (insincerity, bluff, accidental configuration),
rather than to a genuine structural feature of \(S\)'s response.

\textbf{PSC-3 (Programme abandonment).} Upon verification that
\(\Phi_S\) is genuine (via independent observation, simulation, or
repetition), \(O\) abandons the active research, clinical, or
institutional programme organized around \(F_O\) and reorganizes around
\(\Phi_S\). The framework is not augmented or corrected; it is replaced.

\subsubsection{D.2 Gradation and diagnostic
thresholds}\label{d.2-gradation-and-diagnostic-thresholds}

Each marker may be present in strong, attenuated, or absent form.
Diagnostic sensitivity increases when all three are present in strong
form; diagnostic specificity requires all three for confident
attribution to Controfase.

{\def\LTcaptype{none} % do not increment counter
\begin{longtable}[]{@{}
  >{\raggedright\arraybackslash}p{(\linewidth - 2\tabcolsep) * \real{0.5000}}
  >{\raggedright\arraybackslash}p{(\linewidth - 2\tabcolsep) * \real{0.5000}}@{}}
\toprule\noalign{}
\begin{minipage}[b]{\linewidth}\raggedright
Combination
\end{minipage} & \begin{minipage}[b]{\linewidth}\raggedright
Interpretation
\end{minipage} \\
\midrule\noalign{}
\endhead
\bottomrule\noalign{}
\endlastfoot
PSC-1 + PSC-2 + PSC-3 (all strong) & Strong evidence of Controfase
encounter \\
PSC-1 + PSC-2 + PSC-3 (attenuated) & Probable Controfase encounter;
requires further evidence \\
PSC-1 only (strong) & Anomaly within framework; may be refineable; does
not diagnose Controfase \\
PSC-2 only & Conventional artefact hypothesis; does not diagnose
Controfase \\
PSC-3 without PSC-1, PSC-2 & Framework change for other reasons
(institutional, career, etc.); does not diagnose Controfase \\
PSC-1 + PSC-3 without PSC-2 & Rapid framework replacement without
resistance; unusual; requires investigation \\
\end{longtable}
}

\subsubsection{D.3 Application protocols by
domain}\label{d.3-application-protocols-by-domain}

\textbf{Physical-experimental.} Applied to published experimental
reports in the observer's own first-person language (methods, results,
discussion sections) and in interview data with the research team. The
Singh et al.~case is the prototype. Retrospective application to
historical paradigm shifts (the discovery of the neutrino, the 1998
supernova results indicating dark energy, the 2015 gravitational wave
detection) would test generalizability.

\textbf{Clinical.} Applied to the therapist's own reports on client
presentations that resist the therapeutic framework. Distinguishes
Controfase-driven non-response (in which the client's internal
configuration does not actualize the framework's target reactive
circuit) from conventional non-response (in which the framework is
broadly adequate but the specific intervention is ineffective). Full
development of clinical PSC requires methodology design.

\textbf{Institutional.} Applied to institutional reports on systems
(individuals, organizations, sub-units) whose responses do not fit the
institution's governing framework. Distinguishes Controfase-driven
institutional anomalies from conventional non-compliance, resistance, or
dysfunction.

\textbf{Epistemic-theoretical.} Applied to theoretical frameworks
encountering phenomena they do not predict. Distinguishes paradigm-level
disconfirmation (in which the framework fails by type) from
parameter-level disconfirmation (in which the framework fails by
degree).

\subsubsection{D.4 Known limitations}\label{d.4-known-limitations}

\begin{enumerate}
\def\labelenumi{(\alph{enumi})}
\item
  The PSC is observer-side. The same physical or social phenomenon may
  trigger PSC in one observer and not in another, depending on the
  adequacy of each observer's framework. This is a feature, not a bug:
  PSC diagnoses the observer's encounter with Controfase, not Controfase
  as a property of the system absent an observer.
\item
  Strong PSC-1 and PSC-2 in the absence of PSC-3 may indicate Controfase
  encounter followed by epistemic capture (the observer recognizes the
  anomaly but is institutionally or personally unable to reorganize
  around it). This is distinct from absence of Controfase; it indicates
  Controfase detected but not acted upon.
\item
  The PSC can be performed (an observer may perform the three markers
  without having genuinely encountered Controfase) for strategic
  reasons. Distinguishing genuine PSC from performed PSC requires
  triangulation with the observer's subsequent conduct and with
  independent observation of \(S\).
\end{enumerate}

\subsubsection{D.5 Research directions}\label{d.5-research-directions}

\begin{enumerate}
\def\labelenumi{(\roman{enumi})}
\tightlist
\item
  Retrospective PSC analysis of historical paradigm shifts in physics,
  biology, and the cognitive sciences, to test whether the Kuhnian
  revolution concept is co-extensive with Controfase encounter as
  diagnosed by PSC. (ii) Development of a clinical PSC protocol for use
  in therapeutic training, distinguishing client Controfase from client
  dysfunction. (iii) Institutional PSC protocols for analysing
  institutional encounters with anomalous sub-systems. (iv) Integration
  of PSC into the Semantic Genome programme as a criterion for
  observer-recognition of emergent ordinative structures in language
  models.
\end{enumerate}

\begin{center}\rule{0.5\linewidth}{0.5pt}\end{center}

\subsection{References}\label{references}

Aharonov, Y., \& Bohm, D. (1959). Significance of electromagnetic
potentials in the quantum theory. \emph{Physical Review}, 115, 485--491.

Baez, J., \& Stay, M. (2011). Physics, topology, logic and computation:
A Rosetta Stone. In \emph{New Structures for Physics} (pp.~95--172).
Springer.

Bell, J. S. (1964). On the Einstein Podolsky Rosen paradox.
\emph{Physics}, 1, 195--200.

Berry, M., Chambers, R., Large, M., Upstill, C., \& Walmsley, J. (1980).
Wavefront dislocations in the Aharonov-Bohm effect and its water wave
analogue. \emph{European Journal of Physics}, 1, 154.

Coste, C., Lund, F., \& Umeki, M. (1999). Scattering of dislocated wave
fronts by vertical vorticity and the Aharonov-Bohm effect. I. Shallow
water. \emph{Physical Review E}, 60, 4908.

Ghioni, F. (2025). \emph{Ordinative Set Theory: A Concise Guide for AI}
(v2.1). Ordinative Sciences Press.

Ghioni, F. (2025). \emph{OST Extension: Teleodynamics and Causal
Inversion} (v1.1). Ordinative Sciences Press.

Ghioni, F. (2026). \emph{TE\_CORE v5.1}. Technology of Expressions
framework documentation.

Ghioni, F. (2026). \emph{The Collapse Equation: Predicting Phase
Transitions --- Reaction-Diffusion Dynamics of Civilizational Systems}
(v1.2-beta; later versions in preparation). Zenodo.
https://doi.org/10.5281/zenodo.19932311

Ghioni, F. (in preparation). \emph{The Controfase: Universal Operator of
Ordinative Dynamics}. Foundational volume.

Klir, G. J. (1969). \emph{An Approach to General Systems Theory}. Van
Nostrand Reinhold.

Kochen, S., \& Specker, E. P. (1967). The problem of hidden variables in
quantum mechanics. \emph{Journal of Mathematics and Mechanics}, 17,
59--87.

Kuhn, T. S. (1962). \emph{The Structure of Scientific Revolutions}.
University of Chicago Press.

Lakatos, I. (1970). Falsification and the methodology of scientific
research programmes. In I. Lakatos \& A. Musgrave (Eds.),
\emph{Criticism and the Growth of Knowledge} (pp.~91--196). Cambridge
University Press.

Maturana, H., \& Varela, F. (1980). \emph{Autopoiesis and Cognition: The
Realization of the Living}. Reidel.

Okinawa Institute of Science and Technology (OIST). (2026, April 20).
Making waves: water simulation of famous quantum effect reveals
unexpected patterns. \emph{OIST News}. https://www.oist.jp/node/47573
--- source of the research-team statements quoted in §1.2 (Rønning,
Singh, Bandi).

Prigogine, I., \& Stengers, I. (1984). \emph{Order Out of Chaos: Man's
New Dialogue with Nature}. Bantam.

Singh, A., Rønning, J., Liu, C.-C., Angheluta, L., Concha, A., \& Bandi,
M. M. (2026). Topology made visible through standing waves in a spinning
fluid. \emph{Communications Physics}, 9, 123.
https://doi.org/10.1038/s42005-026-02603-w

\begin{center}\rule{0.5\linewidth}{0.5pt}\end{center}

\emph{End of working paper v2.1. Subsequent versions: v2.2 may
incorporate corrections from external reading; v3.0 will follow upon
experimental verification of the extensions of §10 or upon cross-domain
operationalisation of the PSC.}

\end{document}
